\documentclass{article}
\usepackage[russian]{babel}
\usepackage{amsmath}

\begin{document}

Правила редукции для комбинаторов $S$, $K$ и $I$:
\begin{align*}
S x y z &\rightarrow_\beta x z (y z) \\
K x y &\rightarrow_\beta x \\
I x &\rightarrow_\beta x
\end{align*}

На первом шаге применяем комбинатор $S$, на втором — комбинатор $K$:
\begin{align*}
(S K K) x &\rightarrow_\beta K x (K x) \\
&\rightarrow_\beta x
\end{align*}

Так как по определению $I x \rightarrow_\beta x$, мы получаем равенство результатов:
\[ (S K K) x \rightarrow_\beta^* x = I x \]

Аксиома экстенсиональности утверждает, что если две функции возвращают одинаковый результат для любого аргумента $x$ (то есть $F x = G x$), то эти функции тождественны ($F = G$).

Так как $(S K K) x = I x$ для любого $x$, по аксиоме экстенсиональности:
\[ S K K = I \]

\end{document}
